\documentclass[11pt]{article}
\usepackage[margin=1in]{geometry}
\usepackage{booktabs}
\usepackage{hyperref}
\usepackage{longtable}
\usepackage{xcolor}
\usepackage{fancyvrb}

\title{Independent Replication Report:\\
Gopinath et al.\ (2022) --- Phylogenomic Analysis of \emph{Salmonella enterica} Serovar Bovismorbificans}
\author{Ollie (OpenClaw AI) \\ Independent-Replication Project (BVBRC-37, Wave 2026-07-01)}
\date{}

\begin{document}
\maketitle

\begin{abstract}
We independently replicate the core phylogenomic and epidemiological claims of Gopinath et al.\ (\emph{Microorganisms} \textbf{10}(6):1199, 2022) on \emph{Salmonella enterica} subsp.\ \emph{enterica} serovar Bovismorbificans. Working from the paper's own NCBI BioProject (PRJNA378379), we pulled all 82 available Bovismorbificans assemblies, re-confirmed the serovar (82/82 by SeqSero2, antigenic profile \texttt{8:r:1,5}), re-derived the MLST distribution (dominant STs $\{142,377,1499,2640\}$ + minority ST150), and reproduced the paper's headline \textbf{two-polyphyletic-cluster} topology using an independent mash + hierarchical-clustering pipeline. AMR and virulence gene classes were reproduced at the feature-class level via AMRFinderPlus, including the \emph{pVirBov}-associated \emph{spv} marker in 56/82 genomes. The paper-specific 2690-locus core-genome schema and the digital DNA-tiling-array SARA/SARB mining were \emph{not} independently rebuilt (custom, non-redistributable). \textbf{Verdict: REPLICATED}, LLM-judge coverage $\approx 0.92$ (free Argo \texttt{gpt-5.2}).
\end{abstract}

\section{Paper Summary}
Gopinath et al.\ analyze \textbf{95 \emph{S.}\ Bovismorbificans strains} (81 newly whole-genome-sequenced from Switzerland, USA, and Canada --- 69 clinical, 9 food, 1 feed, 1 animal, 1 environment --- plus 14 U.S.\ hummus-outbreak genomes and additional NCBI downloads). They construct a bespoke \textbf{2690-locus core-genome schema} from 150 complete \emph{Salmonella} genomes, apply a \textbf{k-mer-binning} strategy to a broader $>$260-strain public set, and mine a \textbf{digital DNA tiling-array} representation of the legacy SARA/SARB reference collections for near-neighbors. The headline biological result is that \emph{S.}\ Bovismorbificans is \textbf{polyphyletic}, splitting into \textbf{two distinct lineages}: a shared backbone composed of \textbf{ST2640, ST142, ST1499, ST377}, and a separate \textbf{ST150} lineage. They also catalog AMR, prophage, plasmid, and virulence-factor gene content, including the \emph{pVirBov} virulence plasmid and its \emph{spv} genes. The 81 newly-sequenced assemblies are deposited at NCBI BioProject \textbf{PRJNA378379} under the FDA-CFSAN GenomeTrakr umbrella. Open access, CC~BY (DOI: \href{https://doi.org/10.3390/microorganisms10061199}{10.3390/microorganisms10061199}; PMC9228720).

\section{Claims Tested}
\begin{longtable}{clllc}
\toprule
\# & Claim & Type & Public-data testable? & Tested \\
\midrule
C1 & 81 new WGS assemblies public under PRJNA378379 & Data availability & Yes (NCBI Datasets REST) & \checkmark \\
C2 & All strains are serovar Bovismorbificans (\texttt{8:r:1,5}) via SeqSero2 & Genomic (serotyping) & Yes & \checkmark \\
C3 & Whole-genome analysis resolves two distinct polyphyletic clusters & Phylogenomic & Yes (mash proxy) & \checkmark \\
C4 & STs 2640/142/1499/377 = backbone; ST150 = separate lineage & Phylogenomic + MLST & Yes & \checkmark \\
C5 & Mixed clinical+food multi-country sampling within lineages & Metadata & Yes (BioSample attrs) & \checkmark \\
C6 & AMR, virulence-factor, plasmid gene content present & Genomic (gene content) & Yes (AMRFinderPlus) & \checkmark (classes) \\
C7 & Bespoke 2690-locus schema + k-mer binning + DNA-tile SARA/SARB mining & Method-specific pipeline & Not feasibly reproducible & \ding{55} \\
\bottomrule
\end{longtable}

\section{Method}
All heavy compute ran on \textbf{uicgpu} (8$\times$A100, 255 cores); metadata / clustering / figures / judge ran locally.

\begin{enumerate}
\item \textbf{Dataset identification.} Queried NCBI Datasets REST for \texttt{bioproject/PRJNA378379/dataset\_report} (500/page, paginated). Of 425 genomes in the umbrella project, filtered \texttt{organism == "\ldots serovar Bovismorbificans"} $\rightarrow$ \textbf{82 assemblies}. Verified against Table~1 of the paper by matching BioSample IDs (e.g.\ \texttt{SAMN12657228} $=$ strain \texttt{N14\_0646} $=$ WGS project \texttt{WSDC01}).

\item \textbf{Download.} \texttt{datasets download genome accession --inputfile acc.txt --include genome} $\rightarrow$ single 117\,MB zip, 82/82 validated, flattened to one FASTA per accession.

\item \textbf{Serovar confirmation (C2).} \texttt{SeqSero2\_package.py -m k -t 4} in k-mer mode on assembled genomes; parsed predicted serotype and antigenic profile.

\item \textbf{MLST / ST typing (C4, C1).} \texttt{mlst --scheme senterica\_achtman\_2} (pubMLST 7-gene Achtman scheme) on all 82.

\item \textbf{Whole-genome clustering (C3, C4).} \texttt{mash sketch} $+$ \texttt{mash dist} all-vs-all $82\times 82$; SciPy \textbf{average-linkage hierarchical clustering}, 2-cluster cut (\texttt{fcluster maxclust=2}); dendrogram rendered with ST-colored leaves.

\item \textbf{Source/geography metadata (C5).} Per-accession NCBI BioSample attributes (\texttt{isolation\_source}, \texttt{host}, \texttt{geo\_loc\_name}); categorized clinical vs food vs animal/env.

\item \textbf{AMR $+$ virulence content (C6).} \texttt{amrfinder --organism Salmonella --plus} (AMRFinderPlus 3.12.8, DB 2024-07-22.1) on all 82; tallied Element type, class, and gene.

\item \textbf{Verdict (LLM-judge).} Evidence bundle scored by a free Argo model. \texttt{argo:claude-opus-4.8} returned HTTP 502 (known proxy bug) $\rightarrow$ fell back to \textbf{free} \texttt{argo:gpt-5.2} (never paid).
\end{enumerate}

\textbf{Tool versions.} datasets 18.32.0; SeqSero2 1.3.2; mlst 2.35.0 (senterica\_achtman\_2); mash 2.3; AMRFinderPlus 3.12.8 (DB 2024-07-22.1); SciPy 1.18.0.

\section{Results vs.\ Paper}

\subsection{Data Availability (C1) --- \checkmark}
82 Bovismorbificans assemblies present and downloadable from PRJNA378379 (81 draft Contig-level $+$ 1 Chromosome). Fully covers the paper's 81 newly-sequenced set. BioSamples match Table~1.

\subsection{Serovar Confirmation (C2) --- \checkmark exact}
SeqSero2 $\rightarrow$ \textbf{82/82 = ``Bovismorbificans''}, antigenic profile \texttt{8:r:1,5} uniformly. Matches the canonical \emph{S.}\ Bovismorbificans O:H formula and the paper's SeqSero2 confirmation.

\subsection{MLST / ST Distribution (C1, C4) --- \checkmark}
\begin{center}
\begin{tabular}{lrl}
\toprule
ST & This work (n) & Paper's role \\
\midrule
142 & 49 & dominant backbone ST \\
377 & 14 & backbone ST (U.S.\ hummus outbreaks) \\
1499 & 11 & backbone ST \\
2640 & 5 & backbone ST \\
150 & 2 & \textbf{separate lineage} \\
8700 & 1 & minor variant (SLV of backbone) \\
\bottomrule
\end{tabular}
\end{center}
The four dominant STs are exactly $\{142, 377, 1499, 2640\}$ --- precisely the four the paper names as the shared-backbone lineage --- with ST150 present as the minority separate lineage.

\subsection{Two Polyphyletic Clusters (C3, C4) --- \checkmark topology reproduced}
Average-linkage hierarchical clustering of the mash whole-genome distance matrix at 2-cluster cut:

\begin{center}
\begin{tabular}{crl}
\toprule
Cluster & n & Composition \\
\midrule
1 & 2 & ST150 only \\
2 & 80 & ST142 (49) $+$ ST377 (14) $+$ ST1499 (11) $+$ ST2640 (5) $+$ ST8700 (1) \\
\bottomrule
\end{tabular}
\end{center}
This is a direct, independent reproduction of the paper's central claim: \emph{S.}\ Bovismorbificans partitions into two distinct lineages, with the $\{2640,142,1499,377\}$ backbone in one cluster and ST150 forming its own separate lineage. (\texttt{evidence/dendrogram.png}, ST-colored leaves.)

\subsection{Clinical $+$ Food, Multi-Country Sampling (C5) --- \checkmark}
BioSample metadata: \textbf{70 clinical/human, 8 food, $+$ animal/env/feed} (paper: 69 clinical, 9 food, 1 feed, 1 animal, 1 env among the 81). Geography: \textbf{Switzerland 75, Canada 5, USA 2} --- the paper's three collections (U.\ Zurich, Health Canada, U.\ Wisconsin/CFSAN). Clinical and food isolates co-occur within the same STs: ST377 $=$ 7 clinical $+$ 5 food; ST1499 $=$ 9 clinical $+$ 2 food; ST142 $=$ 47 clinical $+$ 1 food.

\subsection{AMR $+$ Virulence Content (C6) --- \checkmark feature classes}
AMRFinderPlus over 82 genomes: \textbf{799 virulence hits, 199 AMR hits, 205 stress/metal hits}.
\begin{itemize}
\item Intrinsic efflux \texttt{mdsA}/\texttt{mdsB} in all 82 (core \emph{Salmonella} multidrug efflux).
\item Acquired AMR is sparse (as expected for this serovar): \texttt{sul2}$\times$5, \texttt{tet(A)}$\times$5, aminoglycoside \texttt{aph}$\times$several, \texttt{blaTEM-1}$\times$2, \texttt{blaCTX-M-55}$\times$1, \texttt{qnrB19}/\texttt{qnrS1}, \texttt{floR}, \texttt{tet(M)}, \texttt{aac(3)-IId}. AMR classes: EFFLUX 164, AMINOGLYCOSIDE 12, SULFONAMIDE 6, TETRACYCLINE 6, BETA-LACTAM 4, QUINOLONE/PHENICOL/BLEOMYCIN/TRIMETHOPRIM few.
\item Universal core virulence: \texttt{invA}, \texttt{avrA}, \texttt{iroB}/\texttt{iroC} (salmochelin), \texttt{sinH}, \texttt{sspH2}, \texttt{sodC1}, \texttt{lpfB}, \texttt{sseK2} in $\approx$all genomes; \texttt{spvD} (\emph{spv} virulence-plasmid marker) in \textbf{56/82} --- directly consistent with the paper's \emph{pVirBov} / virulence-plasmid discussion.
\end{itemize}

\subsection{Not Reproduced (C7) --- \ding{55} honest scope}
The paper's bespoke 2690-locus core-genome schema (built from 150 complete genomes), the k-mer-binning survey of $>$260 strains, and the digital DNA microarray/tiling-array SARA/SARB near-neighbor mining were not reconstructed --- these need the custom schema and legacy array data that are not redistributable as drop-in artifacts. Clustering here used mash genome distance as an independent proxy, which recovers the identical two-lineage topology, so the \emph{conclusion} is replicated even though the exact pipeline is not.

\section{Genuine Critique}
This section is a deliberately unflattering assessment of what our replication does \emph{not} prove, where the evidence is thinner than the verdict tag might suggest, and where we took convenient shortcuts.

\begin{enumerate}
\item \textbf{Mash is a genome-distance proxy, not a phylogeny.} We recovered the paper's two-lineage \emph{topology} at a 2-cluster cut of average-linkage clustering on a mash distance matrix. This is not the same as a bootstrapped ML/parsimony core-genome tree. We did not compute branch supports, did not perform cluster-stability analysis (e.g.\ silhouette / gap statistic), did not vary the linkage method (single/complete/Ward would give somewhat different geometries), and did not compare against an SNP-based tree from a reference-mapped alignment. The two-cluster split ``pops out'' partly because the ST150 signal is so strong; a finer inspection could reveal within-backbone sub-structure the paper's own richer schema would resolve.

\item \textbf{$n=2$ for the ``separate lineage''.} The ST150 cluster is only two genomes wide in our set. Two isolates is a very weak evidentiary base for calling a lineage --- if either of those two assemblies has a large chunk of contamination or an assembly artifact, the separation could be inflated. The paper's own ST150 support draws on broader schema/microarray context we did not reconstruct.

\item \textbf{Serovar call has a known circularity.} SeqSero2 in k-mer mode is highly accurate for Bovismorbificans, but the input genomes were already retrieved by filtering NCBI's organism annotation for ``\ldots serovar Bovismorbificans.'' A 100\% match therefore partly reflects consistency of NCBI's own metadata rather than a blind reconfirmation from raw reads. A stronger test would have started from a mixed \emph{Salmonella} draw and let SeqSero2 select.

\item \textbf{AMR/virulence agreement is at the \emph{class} level.} We reproduced the \emph{feature classes} (intrinsic efflux ubiquitous; acquired AMR sparse; \emph{spv} present in a majority) but did not perform gene-by-gene, isolate-by-isolate concordance against the paper's supplementary tables. It is possible the paper and this replication disagree on individual gene calls in specific strains --- we simply did not check. Similarly, prophage and plasmid replicon typing were not attempted.

\item \textbf{The paper's actual method was not run.} The 2690-locus custom schema, the k-mer binning of $>$260 strains, and the SARA/SARB DNA-tiling-array mining are the paper's genuinely novel methodological contributions. We reproduced only the biological conclusions, not the method. A ``methods-reproducible'' verdict is stronger than a ``conclusions-reproducible'' verdict, and we should be honest we're claiming the weaker one.

\item \textbf{Numerical mismatches are glossed over.} We recovered 82 Bovismorbificans genomes from PRJNA378379 vs.\ the paper's 81 newly-sequenced. We report 70 clinical + 8 food vs.\ paper's 69 clinical + 9 food. These small deltas are almost certainly due to post-publication BioSample metadata edits or a Table-1-vs-BioProject accounting difference, but we did not per-genome reconcile against Table~1.

\item \textbf{Single-judge caveat.} The verdict is scored by a single LLM judge (free Argo \texttt{gpt-5.2}). The intended first-choice judge (\texttt{argo:claude-opus-4.8}) errored (HTTP 502) and we did not retry post-fix. There is no ensemble, no self-consistency vote, and no human adjudication.

\item \textbf{Provenance depth is uneven.} We fully log tool versions and API endpoints, but we do not archive the exact NCBI Datasets snapshot (their contents drift), the exact AMRFinderPlus DB version was pinned but not itself checksum-recorded in the report, and the input FASTAs were not hashed. Full byte-level provenance would require SHA-256 manifests per artifact.

\item \textbf{Coverage estimate ($\approx 0.92$) is a soft number.} It comes from the judge, not from a claim-by-claim scoring rubric we defined a priori. Treat it as ordinal (``high''), not cardinal.
\end{enumerate}

None of the above overturns the \textbf{REPLICATED} verdict --- the biological conclusions were reproduced from independent public data and independent code paths, and no paper claim was contradicted. But the verdict should be read as ``conclusions replicated, method not attempted, evidence strong but not exhaustive.''

\section{Verdict}
\textbf{REPLICATED.} Five of six directly-testable core claims (C1--C5) are independently reproduced on real public data with exact or near-exact agreement, including the paper's headline two-polyphyletic-cluster structure and the precise dominant-ST composition. C6 (AMR/virulence/plasmid features) is reproduced at the feature-class level. Only the paper-specific custom-schema/microarray pipeline (C7) was out of reach, and an independent method (mash clustering) confirms the same biological conclusion. No claim was contradicted. LLM-judge (free Argo \texttt{gpt-5.2}): REPLICATED, coverage $\approx 0.92$.

\section{Open Questions}
A curated set of five genuinely open questions --- not superficial nits, but the follow-up work that would either meaningfully strengthen this replication or unlock new science on \emph{S.}\ Bovismorbificans phylogenomics --- is provided as a structured file: \texttt{open\_questions.json} in this report directory. Each entry carries a \texttt{q} (question), a \texttt{basis} (why it matters, grounded in the paper), and \texttt{next\_steps} (concrete methodology to answer it).

\end{document}
